\section{Determining the kinetic drift}
\label{cal:section}

The two coefficients computed here agree with the first two
sigma-model renormalization coefficients
\cite{BrezinZinnJustin,BrezinZinnJustinLeGuillou,Shin}. We determine them
in the precise lattice normalization rather than importing a change of
renormalization scheme. The bare variable used below is $\beta$;
the fluctuation parameter in Section~\ref{rg:section} is
$g=\beta^{-1/2}$, whereas the usual perturbative charge is $\beta^{-1}$.

The preceding construction proves that the scalar and kinetic
coefficients are well defined and that their canonical parts converge
with the free history. It remains to determine two numbers. We do so by
computing a finite-volume partition function in two ways. The bare
calculation below is an ordinary finite-dimensional Laplace expansion;
its remainder need not be uniform in the volume. Uniformity enters only
in the second calculation, through the coefficient bounds already
proved for the block transformation.

Let $\alpha_h,\kappa_h$ be the canonical kinetic kicks at depth $h$,
starting with zero canonical interactions. Thus a single step has the
form
\[
 \beta'=\beta+\alpha_h+\kappa_h/\beta+\Delta_h.
\]
The error $\Delta_h$ is the regular-error contribution, separately
bounded in Theorem~\ref{thm:rg-closure}. It is absent from the formal
canonical calculation. Define
\[
 B_s=\sum_{h<s}\alpha_h,\qquad C_s=\sum_{h<s}\kappa_h.
\]

\begin{theorem}
\label{thm:calibration}
For the nearest-neighbor O(4) normalization of
\eqref{eq:measure}, and $\gamma=(\log L)/\pi$,
\begin{align}
 B_s&=-\gamma s+O_L(1),&
 C_s&=-\frac{\gamma}{2\pi}s+O_L(1),\label{cal:sums}\\
 \alpha_h&=-\gamma+O_L(\rho^h),&
 \kappa_h&=-\frac{\gamma}{2\pi}+O_L(\rho^h)
 \label{cal:kicks}
\end{align}
for some $\rho<1$. The constants are independent of the number of
subsequent block transformations.
\end{theorem}

\subsection{A finite-torus expansion}

For this calculation only, take spins on $S^D$, so that there are $D$
tangent components. Let $v=n^2$. Write $A(n)$ and $Z_2(n)$ for the
coefficients of $\beta^{-1}$ and $\beta^{-2}$ in $\log Z_\beta(n,n)$ at
fixed $n$. A subscript ${\rm d}$ denotes the operation
$F_{\rm d}(n)=4F(n)-F(2n)$, and put $D_0=3D/2$.

Near an aligned configuration choose the mean-spin axis and write
\[
 q_x=(\sqrt{1-|u_x|^2},u_x),\qquad \sum_xu_x=0,
 \qquad z_x=\sqrt\beta\,u_x.
\]
Integrating over the mean axis gives the chart density
\[
 \prod_x(1-|u_x|^2)^{-1/2}
 \left(v^{-1}\sum_x\sqrt{1-|u_x|^2}\right)^D.
\]
Indeed the delta constraint that the mean tangent component vanishes
has axis Jacobian $|\bar q|^{-D}$, so the displayed factor cancels that
Jacobian. The root with positive mean component is selected. At fixed
volume every configuration outside a neighborhood of the aligned
manifold has a positive action excess, and hence contributes
exponentially little as $\beta\to\infty$.

Let $K$ be the nearest-neighbor Laplacian and let $G_x$ be its mean-zero
inverse on the torus. In this subsection put
\[
 G=G_0,\quad M=v^{-1}\sum_xG_x^2,\quad
 r=\frac{1-v^{-1}}2,\quad c=G-r/2=G_{e_1}.
\]
The components of $z$ are independent centered Gaussian fields with
covariance $G$. Put $R_x=|z_x|^2$. Expansion of the action and chart
density gives
\begin{align}
 L_1={}&\frac12(1-D/v)\sum_xR_x
             -\frac18\sum_{\langle xy\rangle}(R_x-R_y)^2,
 \label{cal:L1}\\
 L_2={}&\left(\frac14-\frac D{8v}\right)\sum_xR_x^2
       -\frac D{8v^2}\left(\sum_xR_x\right)^2
       -\frac1{16}\sum_{\langle xy\rangle}
                  (R_x-R_y)^2(R_x+R_y).
 \label{cal:L2}
\end{align}
All edge sums in these two formulas are unoriented. Therefore
\[
 A(n)=\E L_1,\qquad
 Z_2(n)=\E L_2+\tfrac12\operatorname{Var}L_1.
\]

The first coefficient is
\begin{equation}
\label{cal:A}
 A(n)=\frac{Dv r^2}{4}-\frac{D(D-1)}2G.
\end{equation}
For example $\E(R_x-R_y)^2=4D(G^2-c^2)$ on an edge, which gives
\eqref{cal:A} directly. For the next coefficient, the independent
Gaussian variables $(z_x+z_y)/\sqrt2$ and $(z_x-z_y)/\sqrt2$ have
covariances $(G+c)I$ and $(G-c)I$. Their moments give
\begin{align}
 \E L_2={}&D(D+2)
 \left[\left(\frac v4-\frac D8\right)G^2-vrG^2+\frac{vr^2G}{4}\right]
       -\frac{D^3G^2}{8}-\frac{D^2M}{4}.
 \label{cal:EL2}
\end{align}

To compute the variance, decompose $L_1$ into its second and fourth
Gaussian chaoses, that is, its normal-ordered homogeneous Gaussian
polynomials. The quadratic coefficient matrix of its second chaos is
\[
 -\tfrac12[cK+(D-1)I/v].
\]
This follows by replacing the second-chaos part of
$(R_x-R_y)^2$ by
$4G(:R_x:+:R_y:)-8c:z_x\cdot z_y:$ and summing the four incident edges.
Since $KG$ is the projection off constants, its contribution to half
the variance is
\begin{equation}
\label{cal:chaos2}
 \frac D4\{c^2(v-1)+2c(D-1)G+(D-1)^2M\}.
\end{equation}
The fourth chaos is $-\frac18\sum_{xy}K_{xy}:R_xR_y:$.
In its covariance there are four pairings which send all fields at
$x$ to one endpoint and all fields at $y$ to the other, four with the
endpoints exchanged, and sixteen mixed pairings. The first eight have
color factor $D^2$ and the last sixteen factor $D$. Thus its
contribution to half the variance is
\begin{equation}
\label{cal:chaos4}
 v\left(\frac{D^2T_1}{16}+\frac{DT_2}{8}\right),
\end{equation}
where
\begin{align*}
 T_1&=\sum_x\bigl(KG_\bullet^2\bigr)_x^2,\\
 T_2&=\sum_yK_{0y}\sum_xG_xG_{x-y}
                         \bigl(K(G_\bullet G_{\bullet-y})\bigr)_x.
\end{align*}
Distinct chaoses are orthogonal, so no further variance terms occur.

Here are useful reductions of these expressions. In sums indexed by
$a,d$, let these indices run over all four nearest-neighbor directions.
Set
\[
 J_a(x)=G_{x+a}-G_x,\quad P(x)=\sum_aJ_a(x)^2,
 \quad h_{ad}(x)=J_d(x+a)-J_d(x),
\]
and
\[
 j_0=\sum_xP(x)^2,\qquad
 j_1=\sum_{x,a,d}J_a(x)^2h_{ad}(x)^2.
\]
The discrete product rule
$K(fg)=fKg+gKf-\sum_a(\nabla_af)(\nabla_ag)$ gives
\begin{align}
 \sum_xG_xP(x)&=G^2-M,\label{cal:GP}\\
 T_1&=4G^2(1-v^{-1})-4Gr^2+j_0,\label{cal:T1}\\
 T_2&=8rG^2-2r^2G-2(G^2-M)/v-W,\label{cal:T2}\\
 W&=2r(G^2+c^2)+j_0/2-j_1/4.\label{cal:W}
\end{align}
For clarity, \eqref{cal:W} follows by applying the product rule to
$W=\sum_{x,d}G_xJ_d(x)(K(G_\bullet J_d))_x$.
Its three terms are
\[
 \sum_xG_x(KG)_xP(x),\quad
 \sum_{x,d}G_x^2J_d(x)(KJ_d)_x,
 \quad-\sum_{x,a,d}G_xJ_a(x)J_d(x)h_{ad}(x).
\]
The middle term is $2r(G^2+c^2)$. In the last term use
$2J_dh_{ad}=J_d(x+a)^2-J_d(x)^2-h_{ad}^2$ and sum by parts.
The difference cancels the first term and supplies $j_0/2$;
reversal of each $a$-edge gives
$\sum G_xJ_a h_{ad}^2=-j_1/2$, supplying $-j_1/4$.
The other three identities follow from the same product rule and
$KG=\delta_0-v^{-1}$.

Combining \eqref{cal:EL2}--\eqref{cal:W} yields
\begin{equation}
\label{cal:Z2}
 Z_2(n)=Dv\left[\frac{D-1}{16}j_0+\frac1{32}j_1+\frac1{128}\right]
 -\frac{D(D-1)^2}{4}G^2-\frac{D(D-1)}8G+O(1).
\end{equation}
In fact the remainder displayed as $O(1)$ is exactly
\[
 \frac{D(D-1)(D-2)}4M+\frac{D(D-1)}{8v}G
       -\frac{3D}{128}+\frac{3D}{128v}-\frac D{128v^2}.
\]
This also verifies that the remainder is bounded, using
$G=O(\log n)$ and $M=O(1)$.

\subsection{The logarithm in the finite-size correction}

Let
\[
 G^\infty(x)=\int_{[-\pi,\pi]^2}
 \frac{\cos(p\cdot x)-1}{\sum_\mu4\sin^2(p_\mu/2)}
 \frac{\mathrm d^2p}{(2\pi)^2}
\]
be the infinite-lattice potential, with $G^\infty(0)=0$, and use
superscript $\infty$ for its differences. For $d=1,2,3$ one has,
with $x$ in the centered fundamental square for the torus term,
\begin{align}
 |\nabla^dG(x)|+|\nabla^dG^\infty(x)|&\le C(1+|x|)^{-d},
 \label{cal:green-decay}\\
 |\nabla^d(G-G^\infty)(x)|&\le Cn^{-d}\quad(|x|_\infty\le n/3).
 \label{cal:green-comparison}
\end{align}
The bound for $G^\infty$ holds at every lattice point.
To see these bounds, partition the Fourier multiplier
$[\sum_\mu4\sin^2(p_\mu/2)]^{-1}$ into smooth dyadic annuli of radius
$s$. After $d$ differences, integration by parts bounds each annulus by
$C_Ns^d(1+s|x|)^{-N}$. The corresponding torus sum is its periodization.
For $s\gtrsim n^{-1}$ sum its nonzero translates; for smaller $s$ sum
the infinite-lattice estimates directly. These give
\eqref{cal:green-decay}--\eqref{cal:green-comparison}. The same Fourier
comparison gives $G=(2\pi)^{-1}\log n+O(1)$ and $M=O(1)$.
The leading symbol $ia\cdot p/|p|^2$ gives
\[
 J_a^\infty(x)=-\frac{a\cdot x}{2\pi|x|^2}+O(|x|^{-2}).
\]
One may evaluate it by replacing a cutoff at zero by $e^{-|p|^2}$;
the difference has the stated faster decay, and
$e^{-|p|^2}/|p|^2=\int_1^\infty e^{-t|p|^2}\,\mathrm dt$.

The equation and square symmetry improve the local comparison. At
zero its first difference is $1/(4v)$, its pure centered second
difference is $1/(2v)$, and its mixed centered second difference is zero.
Summing the third-difference bound therefore gives
\begin{equation}
\label{cal:green-quadratic}
 J_a(x)-J_a^\infty(x)=\frac{a\cdot x}{2v}+\frac1{4v}
                  +O((1+|x|)^2/n^3),\qquad |x|_\infty\le n/4.
\end{equation}
Put $c_*=(2\pi)^{-1}$. For $1\le|x|\le n/4$, the preceding formulas
imply
\begin{align*}
 P(x)-P^\infty(x)
   &=-2c_*/v+O((v|x|)^{-1}+|x|/n^3+|x|^2/v^2),\\
 P^\infty(x)&=2c_*^2/|x|^2+O(|x|^{-3}).
\end{align*}
In the difference of the squares the only logarithmic sum is
$-8c_*^3v^{-1}\sum_{1\le|x|\le n/4}|x|^{-2}$.
All the displayed error terms, and the tails outside this range, sum to
$O(v^{-1})$. For $j_1$ the infinite summand decays as $|x|^{-6}$ and
its differentiated comparison errors are summable. Consequently
\begin{equation}
\label{cal:j-asymptotic}
 j_0=j_0^\infty-\frac{2\log n}{\pi^2v}+O(v^{-1}),\qquad
 j_1=j_1^\infty+O(v^{-1}).
\end{equation}

Volume terms cancel in the doubling operation. Substituting
\eqref{cal:j-asymptotic} into \eqref{cal:A}, \eqref{cal:Z2} gives
\begin{align}
 A_{\rm d}(n)&=-\frac{3D(D-1)}{4\pi}\log n+O(1),
 \label{cal:direct1}\\
 Z_{2,\rm d}(n)+\frac{A_{\rm d}(n)^2}{2D_0}
       &=-D_0\frac{D-1}{4\pi^2}\log n+O(1).
 \label{cal:direct2}
\end{align}
There is an instructive cancellation in the second formula. Write
$G'=G_{2n}(0)$. The quadratic Green terms reduce to
$D(D-1)^2(G-G')^2/3$, which is bounded. The cross term in
$A_{\rm d}^2/(2D_0)$ cancels the linear Green term in $Z_{2,\rm d}$.
Thus the only unbounded term is the finite-size correction to $j_0$.

\subsection{Computing the same coefficients after blocking}

We now set $D=3$. Choose a fixed sufficiently large dyadic $m$ and put
$n=mL^s$. Insert $s$ normalized observation kernels at parameters
$\beta_h=\beta+B_h+C_h/\beta$, on both the $n$ and $2n$ tori. Pin one
spin on the last layer to remove global rotation. For each fixed $s$
this is a finite-dimensional integral with the same partition function
up to a coupling-independent constant. Here $\beta$ tends to infinity
at fixed $s$; the calculation does not require an admitted trajectory
in the reference bands $[H_j/2,2H_j]$. Its aligned saddle is unique
after pinning, with strictly positive transverse quadratic form.
All other configurations have a positive action excess: zero
microscopic energy forces alignment, and zero observation energy
propagates that alignment through every layer, including the fallback
branch of the observation.

Successive Gaussian completion and Wick integration compute exactly
the same finite Taylor coefficients as direct Laplace expansion. These
are the canonical prescriptions of Lemma~\ref{rg:canonical-contraction}.
Replacing the next formal parameter by
$\beta+B_{h+1}+C_{h+1}/\beta$ changes the kinetic coefficient only by
$O_s(\beta^{-2})$; its two angular factors put that change beyond the
field order in question. Scalar powers are still retained at their
specified order.

We justify the uniform error needed in comparing coefficients as $s$
varies. At a step whose output period is $\bar m=mL^{s-h-1}$, a
finite-period coefficient differs from its folded bulk value only if
the complete contraction paths travel a distance at least
$c\bar m$. Analytic coefficient norms therefore bound this discrepancy
by $C_Le^{-c\bar m}$ per anchor. A previously introduced discrepancy
of size $\delta$ in the ordinary coefficients grows by at most $C_L$
in one formal step: there are at most four fluctuation orders, and all
the coefficient and covariance sums are absolutely summable. There is
no need to impose an affine cancellation on this discrepancy.
Earlier periods are larger by successive factors $L$, so
\[
 \sum_{r\ge0}C_L^{r+1}e^{-cmL^r}<\infty.
\]
The same sum with any fixed polynomial in $r$ and $mL^r$ is finite.
The period factor pays for the sum over anchors, including scalar
Wick contractions. The depth factor also pays for reexpressing their
inverse-coupling powers in terms of $\beta_s$, since
$|B_s-B_h|\le C_L(s-h)$. Thus both inverse-coupling coefficients of
the propagated winding contributions are bounded uniformly in $s$.

At the final fixed period $m$, the pinned Gaussian and canonical
coefficients are bounded uniformly in $s$. All bulk scalar powers at
earlier stages are proportional to volume and cancel in the doubling
operation. The remaining Gaussian pinning power is
$D_0\log\beta_s$: the powers on periods $m$ and $2m$ are respectively
$-D(m^2-1)/2$ and $-D(4m^2-1)/2$, whose doubling difference is $D_0$.
The blocked expansion consequently has the form
\[
 \begin{gathered}
 \Delta_\beta(mL^s)
 =D_0\log\beta_s+\mathcal C_s
       +\frac{X_s}{\beta_s}+\frac{Y_s}{\beta_s^2}
       +O_s(\beta_s^{-3}),\\
 |X_s|+|Y_s|\le C_{L,m}.
 \end{gathered}
\]
Here $\mathcal C_s$ is independent of $\beta$, and $X_s,Y_s$ include
the final fixed-period coefficients and all propagated winding terms.
Their bounds are uniform in $s$; the displayed remainder is a
fixed-$s$ Laplace remainder. Substituting
$\beta_s=\beta+B_s+C_s/\beta$ and comparing coefficients gives
\begin{align}
 A_{\rm d}(mL^s)&=D_0B_s+X_s,\label{cal:block1}\\
 Z_{2,\rm d}(mL^s)&=D_0(C_s-B_s^2/2)-B_sX_s+Y_s.
 \label{cal:block2}
\end{align}
Equations~\eqref{cal:direct1}, \eqref{cal:block1} prove the first part
of \eqref{cal:sums}. Adding $A_{\rm d}^2/(2D_0)$ in
\eqref{cal:block2} cancels both $B_s^2$ and $B_sX_s$; then
\eqref{cal:direct2} proves its second part.

The triangular contraction of the canonical coefficients, with the
forcing $O(L^{-h})$ from the free kernels, gives exponential convergence
of each kick. The bounded cumulative formulas \eqref{cal:sums}
identify their limits, proving \eqref{cal:kicks} and the theorem.
Notice that no uniform bare Laplace remainder was used.

\subsection{Length as a function of the bare coupling}

For later use we record the consequence for the exact trajectories.

\begin{proposition}
\label{prop:calibration}
For every admitted nearest-neighbor trajectory of length $N$, starting
with zero canonical interactions, whose terminal coupling
belongs to a fixed compact interval $I_H$ about a sufficiently large
$H$, its initial coupling $\beta$ satisfies
\begin{equation}
\label{eq:calibration}
 N\log L=\pi\beta-\tfrac12\log\beta+O_{L,H}(1).
\end{equation}
In particular $L^{-N}$ is bounded above and below by positive fixed
multiples of $\sqrt\beta e^{-\pi\beta}$.
\end{proposition}

\begin{proof}
Index the successive actual couplings by $\beta_h$, $0\le h\le N$.
Theorem~\ref{thm:rg-closure} and \eqref{cal:kicks} give
\[
 \beta_{h+1}=\beta_h-\gamma-\nu/\beta_h+r_h+e_h,
 \quad \nu=\gamma/(2\pi),\quad
 |r_h|\le C_L\beta_h^{-1-\epsilon},\quad
 \sum_h|e_h|\le C_L
\]
for some $\epsilon>0$. Increase $H$ so that $|r_h|\le\gamma/2$.
For
\[
 \Phi(b)=b/\gamma-(\nu/\gamma^2)\log b
\]
Taylor's formula gives
\[
 \Phi(\beta_{h+1})-\Phi(\beta_h)
 =-1+O_L(\beta_h^{-1-\epsilon'})+O_L(|e_h|),
 \qquad \epsilon'=\min(\epsilon,1)>0.
\]
The two terms of order $\beta_h^{-1}$ cancel. The errors are summable
uniformly in $N$. Indeed, summing the raw drift between indices $i<j$
gives $\beta_i-\beta_j\ge\gamma(j-i)/2-C_L$. Thus a unit interval of
couplings can be visited only a bounded number of times: apply this
inequality to its first and last visits. Summation over unit intervals
now bounds $\sum_h\beta_h^{-1-\epsilon'}$.
Telescoping $\Phi$ and using the fixed terminal interval proves
\eqref{eq:calibration}.
\end{proof}


